\section{Experiments}
\label{sec:exp}

\subsection{Dataset and stories}
We evaluate on the course-provided Test-A story set and supplementary course test files. All stories were obtained from the course repository as tagged text files; no stories were written, edited, or re-tagged by us after extraction. Each file uses \texttt{[SCENE-n]} markers and \texttt{[SEP]} delimiters to separate scenes; character entities are tagged in angle brackets (e.g., \texttt{<Lily>}, \texttt{<Ryan>}). Stories were selected to cover four representative categories required for controlled comparison: single human (Stories 01, 02, 14, 19), animal subjects (Story 03, extra bird story), non-human entities (Story 13, robot), and multi-character interaction (Stories 06, 07 with two humans). No story received manual prompt editing or per-case image retouching after automated generation. The full story text for each evaluated story is preserved in the run logs for reproducibility.

\paragraph{Primary showcase (Story-14).}
A single-character rain narrative (\texttt{<Girl>}: runs in the rain; stands under a roof; watches the rain slow).
This story exercises pronoun resolution (``She''), scene-consistency (roof, rain), and IP-Adapter at $\lambda=0.3$.
It is the main single-entity result for the submitted hybrid pipeline.

\paragraph{Controlled comparison matrix.}
Table~\ref{tab:controlled_comparisons} lists representative Test-A stories covering single human, animal pair, robot, and two-human interaction.

\begin{table}[!t]
  \centering
  \small
  \caption{Controlled qualitative comparison matrix.}
  \label{tab:controlled_comparisons}
  \begin{tabular}{p{0.19\columnwidth}p{0.25\columnwidth}p{0.42\columnwidth}}
    \toprule
    \textbf{Story} & \textbf{Routes} & \textbf{Stress test} \\
    \midrule
    19 Student & Single & One human; case-insensitive spec lookup \\
    03 Cat/Dog & Single vs native & Animal subject types; dual animals \\
    13 Robot & Single vs native & Non-human robot identity \\
    06 Jack/Sara & Multi (native) & Two humans; interaction scenes \\
    14 Girl/rain & Single (primary) & Pronouns; scene consistency; $\lambda$ ablation \\
  \end{tabular}
\end{table}

\subsection{Evaluation protocol}
All scores are \textbf{qualitative} on a 1--5 scale (higher is better) across identity, scene faithfulness, subject count, visual quality, artifacts, and automation robustness (Section~\ref{sec:problem}).
Two authors independently rated a subset of runs; disagreements were resolved by discussion.
We report representative panels rather than aggregate means because the course evaluation is relative and human-judged.

\subsection{Ablation studies}
\label{sec:ablations}

\paragraph{Scene-consistency phrases.}
On Story-14 and setting-sensitive stories, we compare generation with and without the scene-consistency suffix assembled from setting-focus and continuity metadata.

\paragraph{Identity conditioning strength.}
On Story-14, we vary $\lambda \in \{0.1, 0.3, 0.7\}$ for IP-Adapter while holding prompts and seeds fixed where possible.

\paragraph{DiT exploratory runs.}
We run PixArt-$\alpha$ with (a) scene-level inference only, (b) story-joint cross-scene blending, and (c) DreamBooth LoRA trained on 4--8 reference images per character, using the same story texts as the SDXL baseline for side-by-side comparison.

\subsection{Implementation environment}
Experiments used local GPU inference (NVIDIA RTX-class, 24\,GB VRAM), PyTorch with float16, and publicly available checkpoints (SDXL-Turbo, IP-Adapter weights, PixArt-$\alpha$, CLIP ViT-B/32).
The LLM planner uses a text API for structured JSON only; all image synthesis is local.
